\chapter{Recommender System}
\label{ch:crs}

The recommender component personalises \gls{emma}'s quiz mode by directing students toward their weakest specialties based on accumulated session performance. \gls{cf} is used because it identifies latent co-weakness patterns across students: a student weak in Anatomy is systematically more likely to be weak in Surgery, even before being tested in Surgery directly. In the absence of real session data, a synthetic dataset of 200 students was constructed with Gaussian-distributed accuracy scores, each student assigned 1-2 weak specialties (mean accuracy $\approx$ 0.30) and stronger performance elsewhere (mean $\approx$ 0.75). Four algorithms from the Surprise library \cite{hug_surprise_2020} were evaluated using 5-fold cross-validation with \gls{loo} evaluation for the ranking metrics: 

\begin{enumerate}
    \item \gls{svd}: A matrix factorisation method that decomposes the student-specialty performance matrix into latent factors.
    \item \gls{nmf}: Non-negative Matrix Factorisation, which also decomposes the performance matrix but with non-negativity constraints.
    \item \gls{knnb}: A neighbourhood-based \gls{cf} algorithm that identifies similar students based on their performance profiles and recommends specialties that similar students are weak in.
    \item \textbf{NormalPredictor}: A baseline algorithm that predicts ratings based on the distribution of observed ratings without considering student similarity or latent factors.
\end{enumerate}

The figure below shows an example student from the synthetic dataset, illustrating the characteristic pattern of two low-scoring weak specialties against a background of stronger performance elsewhere (See Figure \ref{fig:synthetic_students} in Appendix). Additionally, Gaussian noise models were used to create natural variation of scores between specialties within each student. This naturally created a graph similar to a bell-curve with the median at about 77\% (See Figure \ref{fig:student_scores} in Appendix).

\gls{p@k} scores are low across all algorithms (0.051-0.067 at K=5), but this is a ceiling effect rather than a failure: students have only 3 to 4 weak specialties, making perfect Precision@5 impossible by construction. The theoretical maximum P@5 when a student has 3 weak specialties is 0.60. \gls{hr@k} is the more informative metric. \gls{knnb} achieves HR@5 = 0.335 and HR@10 = 0.740, meaning it successfully identifies at least one of a student's weak specialties for 74\% of students at K=10. This is a strong result on a 200-student synthetic dataset and confirms that the \gls{cf} approach is viable for study recommendation. See Figures \ref{fig:crs_algorithms}, \ref{fig:crs_algorithmic_comparison}, \ref{fig:crs_k5_k10}, and \ref{fig:crs_k5_k10_plot} in Appendix for algorithm comparisons.

The champion algorithm is \gls{knnb}, which achieves the lowest \gls{rmse} (0.2208 vs. NormalPredictor's 0.3109, a 29\% reduction) and the highest HitRate at both K values. \gls{knnb} identifies at least one of a student's true weak specialties for 74\% of students at K=10, compared to 48\% for NormalPredictor. \gls{svd} achieves a comparable HR@10 of 54\% but with a higher \gls{rmse} of 0.2358. \gls{knnb}'s neighbourhood-based approach is particularly well suited to this setting: students with similar weakness profiles cluster tightly in the rating space, making the nearest-neighbour approach more effective than the global latent structure recovered by matrix factorisation. See Figure \ref{fig:crs_svd_pca} for PCA projection of \gls{svd}.